% Paper 3 of 7 in The Windstorm Series
\documentclass[12pt,a4paper]{article}

\usepackage{amsmath}
\usepackage{amssymb}
\usepackage{graphicx}
\usepackage{hyperref}
\usepackage{natbib}
\usepackage{booktabs}

\title{Cross-Substrate Convergence and Decomposition of Serial Decoding Throughput}

\author{Grant Lavell Whitmer III\\
The Windstorm Institute, Fort Ann, NY 12827, USA\\
\texttt{grantwhitmer3@gmail.com}}

\date{March 2026}

\begin{document}

\maketitle

\begin{abstract}
This study examines throughput convergence across six domains---genetics, artificial intelligence, human cognition, music, language, and engineered systems---by measuring or calculating effective information per serial decoding event for 31 systems with vocabulary sizes spanning 64,000-fold. A decomposition framework $I_{\mathrm{eff}} = R_M(\varepsilon) + \Delta_s + \xi_i$ partitions throughput into a universal rate-distortion floor, substrate-specific slack, and system-specific residual. Across 31 systems, 94\% of observed throughput values fall between 2 and 6 bits, with a median of 4.39 bits (bootstrap 95\% CI: 3.82--4.67). Monte Carlo sampling (100,000 draws) places 90.5\% of bio-plausible samples in the 3--6 bit band versus 28.3\% for random broad sampling ($p < 10^{-300}$). Three independent evolutionary simulations converge to $K \sim 19$--30, and a co-evolutionary simulation discovers $K = 19.76$ and $\varepsilon = 6.62 \times 10^{-3}$---within 10\% of the ribosome's values---starting from no biological priors. Without AI models, the remaining systems do not cluster significantly tighter than alphabet sizes predict ($p \sim 0.56$), indicating that cross-domain convergence statistically depends on the AI data.
\end{abstract}

\noindent\textbf{Keywords:} basin decomposition; evolutionary simulation; Monte Carlo significance; receiver optimization; cross-domain convergence; alphabet size; discrimination cost; substrate slack; co-evolutionary discovery; bio-plausible sampling

\section{Introduction}

In a companion paper~\cite{whitmer2026rlf}, we established the rate-distortion floor $R_M(\varepsilon)$ as a mechanistic bound on serial decoding throughput and demonstrated its predictive power for the ribosome and its vocabulary-independence for AI tokenizers. This paper asks the broader question: does the same constraint govern serial decoders across all substrates?

We survey 31 systems across six domains, decompose the throughput into floor plus slack plus residual, run extensive computational experiments, and explore the implications for cross-substrate kinship. Evidence is classified into three tiers: [MI] direct mutual information or entropy-rate estimates (Tier~1), [SS] support-size upper bounds computed as $\log_2$ of alphabet size (Tier~2), and [AN] illustrative analogies included for context but excluded from statistical analysis (Tier~3). Only Tier~1 entries are used in formal statistical tests.

\section{Materials and Methods}

\subsection{Throughput Basin Decomposition}

For any serial decoding system, the effective information per event is decomposed as:
\begin{equation}
    I_{\mathrm{eff},i} = R_M(\varepsilon_i) + \Delta_s(i) + \xi_i
\end{equation}
where $R_M(\varepsilon)$ is the rate-distortion floor (universal, mechanistic), $\Delta_s$ is substrate-specific slack (excess capacity margin), and $\xi_i$ is system-specific residual (measurement noise, non-optimality).

\subsection{Data Collection: 31 Systems Across Six Domains}

Throughput values were measured or calculated for 31 serial decoding systems organized into six domains:

\textbf{Biological systems (genetics):} Standard genetic code ($M = 21$), selenocysteine variant ($M = 22$), pyrrolysine variant ($M = 23$).

\textbf{Human phonology:} Cross-linguistic median ($M = 31$), English phonemes ($M = 44$), Hawaiian phonemes ($M = 13$), English consonants from Miller and Nicely confusion matrices~\cite{millernicely1955} ($M = 16$, Tier~1).

\textbf{Human cognition:} Absolute identification ($\sim$7 levels, Tier~1), Hick's Law choice reaction time ($\sim$8 choices, Tier~1), Miller's Law $7 \pm 2$ items~\cite{miller1956}, visual working memory ($\sim$3--4 bits, Tier~1).

\textbf{Music:} Western chromatic scale (12 tones), diatonic major (7 tones), pentatonic (5 tones), Indian classical shrutis (22 tones).

\textbf{Engineered systems:} Braille (64 symbols), semaphore flags (30 symbols), maritime signal flags (40 symbols).

\textbf{AI systems:} GPT-4, LLaMA-3-70B, Claude 3.5 Sonnet, Gemini 1.5 Pro, GPT-2 (model-dependent cross-entropy, designated [MI*]).

\subsection{Monte Carlo Basin Significance Test}

To assess whether the 3--6 bit clustering is statistically significant, 100,000 samples were drawn from each of two distributions:
\begin{itemize}
    \item \textbf{Bio-plausible:} $M$ sampled uniformly from $[4, 64]$, $\varepsilon$ sampled log-uniformly from $[10^{-4}, 0.1]$.
    \item \textbf{Random broad (control):} $M$ sampled uniformly from $[2, 300]$, $\varepsilon$ sampled log-uniformly from $[10^{-5}, 0.5]$.
\end{itemize}

For each sample, $R_M(\varepsilon)$ was computed and basin membership ($R_M \in [3, 6]$ bits) was recorded. Significance was assessed using a binomial test.

\subsection{Evolutionary Simulations}

Three independent evolutionary simulations tested whether cost-constrained optimization converges to biologically relevant alphabet sizes:

\textbf{Implementation A (Genetic Algorithm):} Population of 3,600, 1,000 generations, fitness-based selection.

\textbf{Implementation B (Analytical Cost):} Population of 800, 400 generations, $R_M - \gamma K^\alpha$ cost function.

\textbf{Implementation C (Empirical MI):} Population of 100, 50 generations, empirical mutual information fitness. Each implementation was developed independently to minimize shared assumptions.

A multi-regime substrate sweep varied the cost exponent $\alpha$ from 1.2 (low cost, digital-like) through 1.5 (medium, ribosome-like) to 1.8 (high cost, music-like). A parameter grid sweep covered 3,600 combinations.

\subsection{Co-Evolutionary Simulation}

In the strongest computational test, both $K$ (alphabet size) and $\varepsilon$ (error rate) were allowed to co-evolve under selection pressure with no biological priors, no seeded target, and no knowledge of molecular biology.

\section{Results}

\subsection{Summary Statistics}

Across all 31 systems, the median throughput is 4.39 bits (bootstrap 95\% CI: 3.82--4.67, $n = 10{,}000$ bootstrap resamples). Ninety-four percent of observations fall between 2 and 6 bits. The band $[3, 6]$ contains 71\% of all observations. The histogram shows strong central clustering around 4--5 bits.

Representative values by domain illustrate the convergence: the ribosome reads codons at 4.39 bits per codon; human ears decode English consonants at approximately 4.2 bits per phoneme; the chromatic musical scale discriminates pitches at 3.6 bits per tone; working memory holds items at approximately 3.1 bits each; Morse code transmits at approximately 4.8 bits per symbol; AI transformers process approximately 4.2 bits per token.

\subsection{Substrate Slack}

\begin{table}[ht]
\centering
\caption{Substrate slack across four serial decoding systems.}
\label{tab:slack}
\begin{tabular}{lccc}
\toprule
System & Floor $R_M(\varepsilon)$ & Observed & Slack $\Delta$ \\
\midrule
Ribosome & 4.39 & 4.39 & 0.00 \\
Phonology & 4.71 & 4.95 & 0.24 \\
Chromatic scale & 3.12 & 3.58 & 0.46 \\
Cognition ($M=7$, $\varepsilon=0.12$) & 1.97 & 3.12 & 1.15 \\
\bottomrule
\end{tabular}
\end{table}

Biology operates nearest the floor. Cognition and music carry higher slack, consistent with greater discrimination cost in neural substrates.

\subsection{Monte Carlo Basin Significance}

Bio-plausible parameter sampling places 90.5\% of samples in the $[3, 6]$ bit band versus 28.3\% for random broad sampling (binomial test $p < 10^{-300}$). The enrichment is substantial (62.2 percentage points), demonstrating that the basin is a geometric consequence of biological parameter ranges.

\subsection{Cross-Domain ANOVA}

The Kruskal--Wallis test yields $H = 14.00$, $p = 0.016$. Domains are distinguishable from each other but all fall within the basin. This is expected if substrate-specific costs shift systems within a constrained band.

\subsection{The AI-Dependence Problem}

Without AI models in the dataset, the remaining biological and engineered systems do not cluster significantly tighter than their $M$ values predict (IQR clustering $p \sim 0.56$). The cross-domain convergence story---the claim that ribosomes, phonemes, neurons, and AI transformers converge to the same band---statistically depends on the AI data. The correct interpretation is:

\begin{enumerate}
    \item Within-substrate convergence is real: biological serial decoders cluster in $[3\text{--}6]$ bits because their $M$ and $\varepsilon$ values are constrained by thermodynamics.
    \item Cross-substrate convergence is suggestive but unproven: AI models landing in a similar band may reflect different underlying causes.
    \item The rate-distortion floor is universal: $R_M(\varepsilon)$ applies to any serial decoding system regardless.
\end{enumerate}

\subsection{Evolutionary Simulations}

\begin{table}[ht]
\centering
\caption{Three independent evolutionary simulations.}
\label{tab:evosim}
\begin{tabular}{lccccc}
\toprule
Implementation & Pop & Gens & $\varepsilon$ & Final $K$ & Final bits \\
\midrule
A (GA fitness) & 3,600 & 1,000 & 0.025 & 29.5 & 4.59 \\
B ($R_M - \gamma K^\alpha$) & 800 & 400 & $10^{-4}$ & 19.14 & 4.26 \\
C (Empirical MI) & 100 & 50 & 0.02 & 21.19 & 4.14 \\
\bottomrule
\end{tabular}
\end{table}

Three codebases, three fitness functions, same basin. The multi-regime substrate sweep shows that the cost exponent predicts which biological system the simulation recovers: low cost ($\alpha = 1.2$) yields $K \sim 53$ (complex language), medium cost ($\alpha = 1.5$) yields $K \sim 21$ (ribosome/phonology), and high cost ($\alpha = 1.8$) yields $K \sim 11$ (music/chromatic).

\subsection{Co-Evolutionary Discovery}

When both $K$ and $\varepsilon$ were allowed to co-evolve with no biological priors, the population independently discovered $K = 19.76$ and $\varepsilon = 6.62 \times 10^{-3}$. The ribosome uses $M = 21$ amino acids at $\varepsilon \sim 10^{-4}$ to $5 \times 10^{-3}$. The simulation landed within 10\% of the ribosome's actual values, starting from nothing but math.

This is the paper's most significant computational result: the convergence is not a parameter choice but an emergent property of the optimization landscape that real molecular evolution also found.

\section{Discussion}

\subsection{Mechanistic Explanation}

The decomposition turned the throughput basin from a mysterious convergence into a mechanical explanation: systems cluster in 3--6 bits because the rate-distortion floor constrains the minimum, and substrate-specific costs constrain the maximum. The band is where physics allows serial decoders to operate.

\subsection{The Bifurcated Argument}

Biological systems are trapped in the 3--5 bit band by thermodynamic discrimination costs: Hopfield kinetic proofreading~\cite{hopfield1974} (free energy change $\sim 2$--3 kcal/mol), metabolic constraints ($\sim 10^4$ ATP/bit at synapses), and the Landauer floor. AI systems are trapped in the same band despite having no thermodynamic constraint on vocabulary because excess vocabulary capacity is forcibly reallocated to contextual disambiguation and error correction. Same destination, different prisons.

\subsection{Cross-Substrate Kinship}

The ribosome (3.8 billion years old), the human ear (300 million years old), the transformer attention head (2017), and Braille (1824) all solve the same problem: decode one symbol per time step from a noisy serial stream while minimizing discrimination cost. They are kin---not by ancestry, but by constraint~\cite{whitmer2026fons}.

\subsection{Statistical Caveats}

Between-domain variation is real (Kruskal--Wallis $p = 0.016$). When actual conditional entropy rates are used instead of support-size bounds, literature values shift substantially downward---music to 1.87--2.97 bits/note, phonemes to approximately 3.0 bits, English text to approximately 1.0 bit/character. The convergence band under conditional entropy may be approximately $[1.5, 4.5]$ bits rather than $[3, 6]$. This is an unresolved question that future Tier~1 measurements must address.

\subsection{Thermodynamic Interpretation}

The minimum free energy per discrimination event is $W_{\min} \geq kT \ln(2) \cdot I_{\mathrm{eff}}$. The dimensionless substrate factor $\phi_s = W_s / (kT \ln(2) \cdot I_{\mathrm{eff},s})$ measures how far above the reversible minimum a substrate operates. For biology: $\phi \sim 1.02$. For cognition: $\phi \sim 1.20$. The slack is the thermodynamic price of robustness---cheaper discrimination allows operation closer to the floor, while more expensive discrimination demands a larger margin~\cite{hopfield1974}.

\subsection{Cross-Linguistic Validation}

Coup\'{e} et al.~\cite{coupe2019} measured approximately $39 \pm 5$ bits per second across 17 languages, despite phoneme inventories ranging from 15 to 69. Cross-linguistic data from the World Atlas of Language Structures shows most language families cluster near the median: Indo-European $\sim 4.4$ bits, Austronesian $\sim 3.9$ bits, Niger-Congo $\sim 4.8$ bits. The global mean is 4.56 bits. Only Khoisan languages consistently exceed 6 bits, representing the only systematic outlier from the throughput basin.

\subsection{Suggestive Parallels}

Several systems outside the formal analysis share the $2^6 = 64$ encoding depth: Braille and DNA both use 64 units (different mechanisms---$2 \times 3$ dot grid versus $4^3$ triplets---but the same number); the I Ching (c.\ 1000 BCE) uses 64 hexagrams from 6 binary lines; Base64 encoding uses 64 characters for computationally efficient 6-bit-per-character representation. The recurrence of $2^6 = 64$ across substrates may reflect a deeper optimality of 6-bit encoding depth, though each instance has its own historical explanation and these parallels are noted as suggestive rather than evidentiary.

\subsection{Limitations}

\begin{enumerate}
    \item Metric heterogeneity remains despite tier labeling; different measurement methods may introduce systematic biases.
    \item AI-dependence of cross-domain significance (Section~3.5) means the strongest convergence claim requires further validation.
    \item Between-domain variation is real (Kruskal--Wallis $p = 0.016$), indicating domains are distinguishable even within the basin.
    \item Evolutionary simulations are toy models optimizing designed cost functions, not full physical simulations of molecular evolution.
    \item No Tier~1 measurements (confusion matrices or direct mutual information) exist for most domains; expanding the Tier~1 dataset is the highest priority for future work.
    \item Chinese characters, olfactory receptor throughput, and sign language phonology are not included and represent important gaps in the cross-substrate survey.
    \item The 3--6 bit band is wide relative to the data spread; a more precise basin estimate requires additional Tier~1 measurements.
\end{enumerate}

\section{Conclusion}

The Throughput Basin Decomposition proposes that serial decoding systems share a universal rate-distortion floor plus substrate-specific slack, governed by the cost of reliable discrimination under noise. The observed 3--6 bit clustering across 31 systems is statistically significant against random parameter sampling (90.5\% versus 28.3\%, $p < 10^{-300}$). Three independent evolutionary simulations converge to biologically relevant alphabet sizes ($K \sim 19$--30), and one co-evolutionary simulation independently rediscovered both the genetic code's alphabet size and error rate to within 10\% starting from no biological knowledge. Whether this constrained regime reflects a deep physical law or merely the coincidence of moderate parameter ranges in biological systems remains an open question worthy of future investigation, particularly through expanded Tier~1 confusion-matrix measurements across additional domains.

\vspace{1em}
\noindent\textbf{Author Contributions:} Conceptualization, G.L.W.; methodology, G.L.W.; software, G.L.W.; validation, G.L.W.; formal analysis, G.L.W.; investigation, G.L.W.; resources, G.L.W.; data curation, G.L.W.; writing---original draft preparation, G.L.W.; writing---review and editing, G.L.W.; visualization, G.L.W.; supervision, G.L.W.; project administration, G.L.W.

\noindent\textbf{Funding:} This research received no external funding. All work was self-funded by the author.

\noindent\textbf{Data Availability Statement:} All data, analysis scripts, evolutionary simulation code, and experiment protocols supporting the findings of this study are publicly available at \url{https://github.com/Windstorm-Labs/throughput-basin} (accessed 12 April 2026).

\noindent\textbf{Acknowledgments:} This paper was developed through adversarial review by six frontier AI models (Claude, GPT, Grok, Gemini, Sonar, and Perplexity). Mathematical derivations, data analysis, and evolutionary simulations were performed with the assistance of Claude (Anthropic) and other AI research tools. The tokenizer-sweep experiment (1,749 models) was executed by Hermes OC1 on an NVIDIA RTX 5090. The evolutionary simulations were conducted during adversarial review by Grok and Sonar.

\noindent\textbf{Conflicts of Interest:} The author declares no conflict of interest.

\begin{thebibliography}{14}

\bibitem{whitmer2026rlf}
Whitmer III, G.L.
The Receiver-Limited Floor: Rate-Distortion Bounds on Serial Decoding Throughput.
\textit{Zenodo} \textbf{2026}.
\href{https://doi.org/10.5281/zenodo.19322973}{DOI: 10.5281/zenodo.19322973}

\bibitem{millernicely1955}
Miller, G.A.; Nicely, P.E.
An Analysis of Perceptual Confusions among Some English Consonants.
\textit{J. Acoust. Soc. Am.} \textbf{1955}, \textit{27}, 338--352.
\href{https://doi.org/10.1121/1.1907526}{DOI: 10.1121/1.1907526}

\bibitem{miller1956}
Miller, G.A.
The Magical Number Seven, Plus or Minus Two: Some Limits on Our Capacity for Processing Information.
\textit{Psychol. Rev.} \textbf{1956}, \textit{63}, 81--97.
\href{https://doi.org/10.1037/h0043158}{DOI: 10.1037/h0043158}

\bibitem{hopfield1974}
Hopfield, J.J.
Kinetic Proofreading: A New Mechanism for Reducing Errors in Biosynthetic Processes Requiring High Specificity.
\textit{Proc. Natl. Acad. Sci. USA} \textbf{1974}, \textit{71}, 4135--4139.
\href{https://doi.org/10.1073/pnas.71.10.4135}{DOI: 10.1073/pnas.71.10.4135}

\bibitem{whitmer2026fons}
Whitmer III, G.L.
The Fons Constraint: Information-Theoretic Convergence on Encoding Depth in Self-Replicating Systems.
\textit{Zenodo} \textbf{2026}.
\href{https://doi.org/10.5281/zenodo.19274048}{DOI: 10.5281/zenodo.19274048}

\bibitem{shannon1948}
Shannon, C.E.
A Mathematical Theory of Communication.
\textit{Bell Syst. Tech. J.} \textbf{1948}, \textit{27}, 379--423.
\href{https://doi.org/10.1002/j.1538-7305.1948.tb01338.x}{DOI: 10.1002/j.1538-7305.1948.tb01338.x}

\bibitem{tlusty2008}
Tlusty, T.
Rate-Distortion Scenario for the Emergence and Evolution of Noisy Molecular Codes.
\textit{Phys. Rev. Lett.} \textbf{2008}, \textit{100}, 048101.
\href{https://doi.org/10.1103/PhysRevLett.100.048101}{DOI: 10.1103/PhysRevLett.100.048101}

\bibitem{coupe2019}
Coup\'{e}, C.; Oh, Y.M.; Dediu, D.; Pellegrino, F.
Different Languages, Similar Encoding Efficiency.
\textit{Sci. Adv.} \textbf{2019}, \textit{5}, eaaw2594.
\href{https://doi.org/10.1126/sciadv.aaw2594}{DOI: 10.1126/sciadv.aaw2594}

\bibitem{berger1971}
Berger, T.
\textit{Rate Distortion Theory};
Prentice-Hall: Englewood Cliffs, NJ, USA, 1971.

\bibitem{tishby1999}
Tishby, N.; Pereira, F.C.; Bialek, W.
The Information Bottleneck Method.
In Proceedings of the 37th Allerton Conference on Communication, Control, and Computing, Monticello, IL, USA, 22--24 September 1999; pp.\ 368--377.

\bibitem{zaher2009}
Zaher, H.S.; Green, R.
Quality Control by the Ribosome Following Peptide Bond Formation.
\textit{Nature} \textbf{2009}, \textit{457}, 161--166.
\href{https://doi.org/10.1038/nature07582}{DOI: 10.1038/nature07582}

\bibitem{whitmer2026tb_code}
Whitmer III, G.L.
Throughput Basin: Experiment Code and Data.
\textit{GitHub} \textbf{2026}.
Available online: \url{https://github.com/Windstorm-Labs/throughput-basin} (accessed 12 April 2026).

\bibitem{whitmer2026sdb}
Whitmer III, G.L.
The Serial Decoding Basin: Five Experiments on Convergence, Thermodynamic Anchoring, and Receiver-Limited Geometry.
\textit{Zenodo} \textbf{2026}.
\href{https://doi.org/10.5281/zenodo.19323423}{DOI: 10.5281/zenodo.19323423}

\bibitem{borst1999}
Borst, A.; Theunissen, F.E.
Information Theory and Neural Coding.
\textit{Nat. Neurosci.} \textbf{1999}, \textit{2}, 947--957.
\href{https://doi.org/10.1038/14731}{DOI: 10.1038/14731}

\end{thebibliography}

\vspace{2em}
\noindent\textit{This paper is Paper 3 of The Windstorm Series (Papers 1--7) on universal constraints in serial information processing.}

\vspace{1em}
\noindent\textcopyright~2026 Grant Lavell Whitmer III. This work is licensed under a \href{https://creativecommons.org/licenses/by/4.0/}{Creative Commons Attribution 4.0 International License} (CC BY 4.0).

\end{document}
